% ==============================================================================
% ASTROPHYSICS COURSE COMPANION - GLOBAL PREAMBLE
% ==============================================================================

% ---------- Geometry & Page Setup ----------
\usepackage[a4paper,margin=2.6cm,headheight=14pt]{geometry}
\usepackage[T1]{fontenc}
\usepackage[utf8]{inputenc}
\usepackage{newpxtext}
\usepackage{graphicx}
\graphicspath{{figs/}}
\usepackage{float}
\usepackage{placeins}
\usepackage{comment}

% ---------- Float placement: let text share a page with a figure whenever
% ---------- there is room, instead of pushing the figure to its own page
\renewcommand{\topfraction}{0.9}
\renewcommand{\bottomfraction}{0.8}
\renewcommand{\textfraction}{0.07}
\renewcommand{\floatpagefraction}{0.8}
\setcounter{topnumber}{3}
\setcounter{bottomnumber}{3}
\setcounter{totalnumber}{5}

% ---------- Mathematics & Physics ----------
\usepackage{amsmath}
\usepackage{amssymb}
\usepackage{mathtools}
\usepackage{newpxmath}
\usepackage{microtype}
\usepackage{physics}
\newcommand{\bmmax}{0}\newcommand{\hmmax}{0}  % newpxmath uses many math alphabets
\usepackage{bm}
\usepackage{nicefrac}
\usepackage{dsfont}
\newcommand{\Id}{\mathds{1}}   % identity matrix

% ---------- Fonts, Typography & Formatting ----------
\usepackage{booktabs}
\usepackage{array}
\usepackage{enumitem}
\usepackage{parskip}
\usepackage{listings}
\lstset{
  basicstyle=\small\ttfamily,
  backgroundcolor=\color{panel},
  frame=single,
  rulecolor=\color{linecol},
  keywordstyle=\color{accent2}\bfseries,
  commentstyle=\color{muted}\itshape,
  stringstyle=\color{accent},
  breaklines=true,
  showstringspaces=false
}
\setlength{\parindent}{0pt}
\setlength{\parskip}{0.45em plus 0.1em minus 0.1em}
\raggedbottom
\usepackage{wasysym}

% ---------- Colors ----------
\usepackage{xcolor}
\definecolor{colone}{RGB}{180,60,60}
\definecolor{coltwo}{RGB}{60,110,180}
\definecolor{colthree}{RGB}{60,140,80}

\definecolor{accent}{HTML}{1F4E79}   % headings, labels, links (house colour)
\definecolor{accent2}{HTML}{2B5F6B}
\definecolor{muted}{HTML}{5A564F}
\definecolor{panel}{HTML}{F8F6F0}
\definecolor{linecol}{HTML}{D8D2C4}
\definecolor{ink}{HTML}{1E1C1A}
\definecolor{astroblue}{HTML}{1B4F72}
\definecolor{astrodark}{HTML}{17202A}
\definecolor{sunyellow}{HTML}{E8A33D}

% ---------- Captions, headings and running heads (house style of the notes template) ----------
\usepackage{caption}
\captionsetup{font=small,labelfont={bf,color=accent},margin=1em}
\usepackage{titlesec}
\titleformat{\chapter}[display]{\normalfont\sffamily\huge\bfseries\color{accent}}
  {\large\chaptertitlename\ \thechapter}{8pt}{\Huge}
\titleformat{\section}{\normalfont\sffamily\Large\bfseries\color{accent}}{\thesection}{0.8em}{}
\titleformat{\subsection}{\normalfont\sffamily\large\bfseries}{\thesubsection}{0.8em}{}
\titleformat{\subsubsection}{\normalfont\sffamily\normalsize\bfseries}{\thesubsubsection}{0.8em}{}
\setcounter{tocdepth}{1}
\usepackage{fancyhdr}
\pagestyle{fancy}
\fancyhf{}
\fancyhead[LE,RO]{\small\thepage}
\fancyhead[RE]{\small\sffamily\nouppercase{\leftmark}}
\fancyhead[LO]{\small\sffamily\nouppercase{\rightmark}}
\renewcommand{\headrulewidth}{0.3pt}
\fancypagestyle{plain}{\fancyhf{}\fancyfoot[C]{\small\thepage}\renewcommand{\headrulewidth}{0pt}}

% ---------- Hyperlinks ----------
\usepackage[colorlinks=true,
            linkcolor=accent,
            citecolor=accent,
            urlcolor=accent,
            bookmarksdepth=3]{hyperref}

% ---------- TikZ & Diagramming ----------
\usepackage{tikz}
\usetikzlibrary{calc,arrows.meta,positioning,decorations.pathreplacing,decorations.markings,angles,quotes,shapes.geometric}
\usepackage{tikz-3dplot}

% ---------- Theorem & Definition Environments ----------
\newtheorem{definition}{Definition}[chapter]
\newtheorem{result}{Result}[chapter]

% ---------- tcolorbox Styling ----------
\usepackage[most]{tcolorbox}
% Worked example: a calm grey box with a numbered title, as in the notes template
\newcounter{wex}[chapter]
\renewcommand{\thewex}{\thechapter.\arabic{wex}}
\newtcolorbox{example}[1][]{enhanced,breakable,sharp corners=south,arc=2mm,
  boxrule=0.4pt,left=2.5mm,right=2.5mm,top=1.5mm,bottom=1.5mm,
  fonttitle=\bfseries\sffamily\small,coltitle=black!80,colbacktitle=black!5,
  titlerule=0.3pt,before skip=10pt,after skip=10pt,colback=white,colframe=black!28,
  step=wex,title={Worked example~\thewex\ifx\relax#1\relax\else: #1\fi}}
\tcbset{
  basestyle/.style={
    enhanced, breakable, colback=panel, colframe=linecol,
    boxrule=0.5pt, leftrule=2.5pt, arc=2pt, outer arc=2pt,
    left=10pt, right=10pt, top=7pt, bottom=7pt,
    fonttitle=\bfseries\sffamily\footnotesize,
    coltitle=muted, before skip=10pt, after skip=12pt
  },
  notebox/.style={basestyle, colframe=accent2, colback=panel, coltitle=white, colbacktitle=accent2},
  flagbox/.style={basestyle, colframe=accent, colback=panel, coltitle=white, colbacktitle=accent},
  keybox/.style={basestyle, colframe=astroblue, colback=white, coltitle=white, colbacktitle=astroblue}
}

\newenvironment{keyresult}
  {\begin{result}\begin{tcolorbox}[keybox]}
  {\end{tcolorbox}\end{result}}

% ---------- Labelled equality chains (STYLE_GUIDE rule 11) ----------
% In a chain write  &\eqby{1} ...  and right after the chain explain each label in
% \begin{explain} \item ... \end{explain}  (items numbered (1), (2), ... automatically).
\newcommand{\eqby}[1]{\overset{\mathclap{\text{\scriptsize\color{accent2}(#1)}}}{=}}
\newcommand{\relby}[2]{\overset{\mathclap{\text{\scriptsize\color{accent2}(#1)}}}{#2}}
\newlist{explain}{enumerate}{1}
\setlist[explain]{label={\color{accent2}\small(\arabic*)},leftmargin=2.2em,
  itemsep=1pt,topsep=3pt,parsep=0pt}

% ---------- Working definitions (STYLE_GUIDE rule 31) ----------
% \unpack{...} is a short green bracket after a hard term, used ONLY inside workdef boxes.
\newcommand{\unpack}[1]{{\color{colthree!80!black}\,[\,#1\,]}}
\newcounter{wdef}[chapter]
\renewcommand{\thewdef}{\thechapter.\arabic{wdef}}
\newtcolorbox{workdef}[1][]{notebox,
  title={Working definition \refstepcounter{wdef}\thewdef\ifx\relax#1\relax\else\ (#1)\fi}}

% ---------- Custom Astronomical & Vector Macros ----------
\newcommand{\vecr}{\vec r}
\newcommand{\vecL}{\vec L}
\newcommand{\vecA}{\vec A}
\newcommand{\vecf}{\vec f}
\newcommand{\vecF}{\vec F}
\newcommand{\vecp}{\vec p}
\newcommand{\avg}[1]{\left\langle #1\right\rangle}
\newcommand{\vecx}{\vec x}
\newcommand{\hatn}{\hat n}   % unit vector to the ascending node
\newcommand{\rhat}{\hat r}
\newcommand{\phihat}{\hat\phi}
\newcommand{\thetahat}{\hat\theta}
\newcommand{\note}[1]{\textcolor{blue!60!black}{\textit{[#1]}}}

% Astronomical acronyms
\newcommand{\LST}{\mathrm{LST}}
\newcommand{\GST}{\mathrm{GST}}
\newcommand{\GMST}{\mathrm{GMST}}
\newcommand{\GAST}{\mathrm{GAST}}
\newcommand{\LMST}{\mathrm{LMST}}
\newcommand{\LAST}{\mathrm{LAST}}
\newcommand{\UT}{\mathrm{UT}}
\newcommand{\UTC}{\mathrm{UTC}}
\newcommand{\TAI}{\mathrm{TAI}}
\newcommand{\TT}{\mathrm{TT}}
\newcommand{\TDB}{\mathrm{TDB}}
\newcommand{\JD}{\mathrm{JD}}
\newcommand{\MJD}{\mathrm{MJD}}
\newcommand{\EoT}{\mathrm{EoT}}
\newcommand{\SNR}{\mathrm{SNR}}
\newcommand{\RON}{\mathrm{RON}}
